\scanpage{33}
partielles i.e. l'inclusion des parties de $\Phi$.

Passant au quotient par $\underline a$, on trouve le \uline{cube combinatoire \inserted{\uline{gauche} (ou} projectif)} $\widetilde C$, et l'\uline{octaèdre comb. \inserted{\uline{gauche} (ou} projectif)} $\widetilde O$, associé à (1),
\[
\widetilde C_i=C/\{1,\underline a\}\qquad
\widetilde O_i=O/\{1,\underline a\}\tag*{(4)}
\]
(relations d'incidence obtenues en passant au quotient). L'\deleted{\unsure{es}}octaèdre \deleted{projectif} \inserted{gauche} a donc $3$ sommets $6$ arêtes et $4$ faces -- en fait c'est la décomposition du plan projectif réel par un système de trois droites indexées par $I$. (À tout élément $i\in I$, on fait correspondre la ``droite'' correspondant aux deux ``arêtes'' projectives, qui correspondent aux quatre arêtes ordinaires données par $A\in\mathfrak P_2(\Phi)$,
\figure{\begin{tikzpicture}[x=.009cm,y=-.009cm]
\draw (8,233) .. controls (140,246) and (340,213) .. (528,197);
\draw (135,6) .. controls (239,129) and (399,223) .. (550,389);
\draw (49,413) .. controls (101,343) and (207,260) .. (334,62);
\draw (519,191)--(528,197)--(519,201);
\draw (102,341)--(105,349)--(117,340);
\end{tikzpicture}}
$A\overset{\sim}{\longleftrightarrow}I\setminus\{i\}$.

On ne récupère \uline{pas} le cube ou octaèdre ordinaire, à isom canonique près, à partir de la figure \inserted{``gauche''} ``projective'' correspondante, vu que les groupes d'\unsure{inerties|inertie}, dans le cas ordinaire et dans le cas \deleted{projectif}, \inserted{gauche}, ne sont pas isomorphes, plus précisément, $C\mapsto\widetilde C$ resp. $O\vdash\widetilde{C}$
% archive sheet number: 32
